\section{Introduction}

Bus holding control aims to reduce headway instability by delaying selected vehicles at stops.
The control problem is local in execution but system-wide in consequence: a hold action changes forward and backward headways, passenger waiting, dwell, and downstream bunching.
Modern reinforcement learning (RL) methods can learn holding policies in simulation, but practical deployment is constrained by two persistent gaps.
First, high-fidelity simulation and manual calibration are expensive to reproduce in each city.
Second, public data availability differs sharply across agencies: one city may expose speed bands and passenger OD data, while another may only provide GTFS schedules and coarse ridership summaries.

Hybrid offline-to-online methods such as \Htwo learn from both logged data and simulator interaction by correcting or reweighting dynamics gaps \citep{niu2022h2o,liu2024h2oplus}.
This is a natural starting point for transit systems, where a historical buffer and a lower-fidelity simulator are often available.
The difficulty is cross-city transfer.
A dense monolithic residual model can fit simulator error in the source cities, but its parameters mix mechanisms that may shift differently across cities: schedule headways, route speeds, stop demand, and dwell behavior do not all change together.
If cross-city transfer depends on localized changes in mechanisms, a factored representation should be more stable than a dense one \citep{bengio2020meta,feng2022factored}.

This paper studies that hypothesis in a deliberately conservative setting.
We do not claim that the current data are vehicle-level real-world trajectories.
Instead, we construct full-route open transit city bundles and evaluate models on static-data-derived offline dynamics targets.
The protocol is meant to answer a narrower question: when only open static data and an uncalibrated simulator-like transition generator are available, does a mechanism-factored residual adapter transfer better across cities than a dense \Htwo-style residual correction?
We separate three deployment regimes that are often conflated: single-source zero-shot transfer, multi-source target-label-free transfer, and few-shot target offline adaptation with a small amount of passive target-city history.
This is the bus-holding core of the paper.
Because passive bus logs cannot reveal what would have happened under unobserved hold actions, we use traffic-signal control as an additional counterfactual stress test where SUMO can roll out alternative actions under a shared benchmark protocol.
The traffic-signal experiments are not presented as the same task or metric scale as bus holding; they test whether the same transfer principle can remain safe when strong classical control rules and external RL leaderboard baselines are available.

We introduce Causal-Factored Cross-City Model Transfer (\ours).
\ours fits residual corrections at the level of output mechanisms rather than one dense residual over all outputs.
For each target quantity, it uses a mechanism-specific sparse parent set: headway and gap outputs depend primarily on neighboring headways, target headway, speed, time, and action; demand and dwell outputs depend on waiting passengers, speed, target headway, time, and action; speed outputs use speed, time, and route-position features; reward uses policy-relevant headway and waiting quantities.
We also add a target-label-free source-selection layer.
It first chooses the target-summary feature set by leave-one-source validation, then uses the selected unlabeled target summary to decide when a source-city ensemble is credible; target transition labels are held out for evaluation.
The theoretical motivation is a mechanism-risk decomposition: squared transition error decomposes over output mechanisms, source ensembles obey a convex risk bound under squared loss, and source-only selector quality depends on how well leave-one-source scores approximate target risk rankings.

\begin{figure}[!t]
    \centering
    \includegraphics[width=0.98\textwidth]{causal_transfer_comparison.png}
    \caption{Conceptual comparison of cross-city residual transfer strategies for bus holding. A dense \Htwo-style residual adapter learns one monolithic correction; MC-MC-style latent world-model adaptation represents domain shift implicitly; \ours{} instead transfers auditable headway/gap, demand/dwell, and speed mechanisms through a target-label-free source gate.}
    \label{fig:causal-transfer}
\end{figure}

The paper contributes:
\begin{enumerate}[leftmargin=1.5em]
    \item A cross-city open-transit validation protocol using four full-route city bundles and a strict leave-one-city-out evaluation unit.
    \item A mechanism-wise residual transfer model with source-only target-summary feature selection and a guarded source-city ensemble.
    \item A formal transfer-risk motivation showing how mechanism shift, source weighting, and source-only selector error enter the target risk.
    \item A target-information ladder spanning single-source zero-shot transfer, multi-source target-label-free transfer, and few-shot target offline adaptation.
    \item A counterfactual traffic-signal stress test showing that a pressure-guarded mechanism-transfer controller reaches pressure-rule-level safety on short RESCO rollouts, while long-horizon results favor a classical phase-pressure rule.
    \item Full-budget native-protocol RESCO and LibSignal RL cross-checks, kept separate from the same-protocol \ours phase-MPC comparisons.
    \item An explicit boundary on the claim: the evidence supports target-label-efficient cross-city residual transfer and safe-prior gating, not field deployment without counterfactual or vehicle-level validation.
\end{enumerate}
